% Options for packages loaded elsewhere
\PassOptionsToPackage{unicode}{hyperref}
\PassOptionsToPackage{hyphens}{url}
\PassOptionsToPackage{space}{xeCJK}
\documentclass[
  a4paper,11pt]{ltjsarticle}
\usepackage{xcolor}
\usepackage[margin=1in]{geometry}
\usepackage{amsmath,amssymb}
\usepackage{graphicx}
\setcounter{secnumdepth}{-\maxdimen} % remove section numbering
\usepackage{iftex}
\ifPDFTeX
  \usepackage[T1]{fontenc}
  \usepackage[utf8]{inputenc}
  \usepackage{textcomp}
\else
  \usepackage{unicode-math}
  \defaultfontfeatures{Scale=MatchLowercase}
  \defaultfontfeatures[\rmfamily]{Ligatures=TeX,Scale=1}
\fi
\usepackage{lmodern}
\ifPDFTeX\else
  \setmainfont[]{Times New Roman}
  \ifXeTeX
    \usepackage{xeCJK}
    \setCJKmainfont[]{Hiragino Mincho ProN}
  \fi
  \ifLuaTeX
    \usepackage[]{luatexja-fontspec}
    \setmainjfont[]{Hiragino Mincho ProN}
  \fi
\fi
\IfFileExists{upquote.sty}{\usepackage{upquote}}{}
\IfFileExists{microtype.sty}{%
  \usepackage[]{microtype}
  \UseMicrotypeSet[protrusion]{basicmath}
}{}
\makeatletter
\@ifundefined{KOMAClassName}{%
  \IfFileExists{parskip.sty}{%
    \usepackage{parskip}
  }{%
    \setlength{\parindent}{0pt}
    \setlength{\parskip}{6pt plus 2pt minus 1pt}}
}{% if KOMA class
  \KOMAoptions{parskip=half}}
\makeatother
\usepackage{longtable,booktabs,array}
\newcounter{none}
\usepackage{calc}
\usepackage{etoolbox}
\makeatletter
\patchcmd\longtable{\par}{\if@noskipsec\mbox{}\fi\par}{}{}
\makeatother
\IfFileExists{footnotehyper.sty}{\usepackage{footnotehyper}}{\usepackage{footnote}}
\makesavenoteenv{longtable}
\usepackage{bookmark}
\IfFileExists{xurl.sty}{\usepackage{xurl}}{}
\urlstyle{same}
\usepackage{hyperref}
\hypersetup{
  hidelinks,
  pdfcreator={LaTeX via pandoc}}
\setlength{\emergencystretch}{3em}
\providecommand{\tightlist}{%
  \setlength{\itemsep}{0pt}\setlength{\parskip}{0pt}}
\usepackage{amsthm}
\newtheorem{theorem}{定理}[section]
\newtheorem{proposition}[theorem]{命題}
\newtheorem{lemma}[theorem]{補題}
\newtheorem{corollary}[theorem]{系}
\theoremstyle{definition}
\newtheorem{definition}[theorem]{定義}
\theoremstyle{remark}
\newtheorem{remark}[theorem]{注意}
\newtheorem{observation}[theorem]{観察}
\begin{document}
\section{閉じた球面上の形不変定常波------次元ごとの安定性と3+1時空の幾何学的必然性}\label{ux9589ux3058ux305fux7403ux9762ux4e0aux306eux5f62ux4e0dux5909ux5b9aux5e38ux6ce2ux6b21ux5143ux3054ux3068ux306eux5b89ux5b9aux6027ux306831ux6642ux7a7aux306eux5e7eux4f55ux5b66ux7684ux5fc5ux7136ux6027}

\textbf{著者}: 木原 範昭 (Noriaki Kihara)\\
\textbf{所属}: WF System Co., Ltd.\\
\textbf{ORCID}: 0009-0004-6753-4020\\
\textbf{版}: v1.0\\
\textbf{日付}: 2026 年 4 月

\begin{center}\rule{0.5\linewidth}{0.5pt}\end{center}

\subsection{要旨}\label{ux8981ux65e8}

閉じた \(n\) 次元球面 \(S^n\)
上の位相波動方程式を、1次元から4次元まで系統的に分析する。位相波動方程式は空間曲率
\(R\) に依存せず、\(R\)
が関与するのは定常波条件（大域的境界条件）のみであることを示す。形不変定常波のパラメータ記述
\((\theta_0, \lambda_{\text{base}}, A, N)\)
が全次元で保存されること、保存量 \(\int |u|^2 \, dV\)
が波束の同一性を担うこと、および Huygens
の原理が奇数次元球面でのみ成立し偶数次元で破れることから、\textbf{形不変定常波が安定に存在できる最小の非自明な球面は
\(S^3\)（3次元）}であることを示す。さらに Frobenius
の定理による離散代数側からの独立な論証を加え、連続側と離散側が共に
\textbf{3+1 時空構造} を必然化することを論じる。

\textbf{キーワード}: 位相波動方程式, 形不変定常波, 閉球面, Huygens
の原理, Frobenius の定理, 3+1時空

\begin{center}\rule{0.5\linewidth}{0.5pt}\end{center}

\subsection{§1
導入と主張範囲}\label{ux5c0eux5165ux3068ux4e3bux5f35ux7bc4ux56f2}

本論文は、閉じた \(n\) 次元球面 \(S^n\)
上で位相波動方程式の定常波がどのような次元依存性を持つかを、純粋に幾何学的・数論的な事実として記述する。

本論文の主張は以下の3点に限定される:

\begin{enumerate}
\def\labelenumi{\arabic{enumi}.}
\tightlist
\item
  位相波動方程式は空間曲率 \(R\) と本質的に直交する演算であり、\(R\)
  が関与するのは定常波条件（閉空間の周長と波長の整合性）のみである。
\item
  形不変定常波が波束としての同一性を保って安定に存在できるのは、Huygens
  の原理が成立する奇数次元球面に限られる。最小の非自明な次元は \(S^3\)
  である。
\item
  Frobenius の定理（結合的除算代数が
  \(\mathbb{R}, \mathbb{C}, \mathbb{H}\)
  に限られる）により、離散代数側からも独立に4次元が必然化され、四元数の
  \(1+3\) 分解が \(3+1\) 時空構造を代数的に内在する。
\end{enumerate}

以下の要素は \textbf{本論文の主張範囲外} である:

\begin{itemize}
\tightlist
\item
  力学的安定性（時間発展に対する固有値解析、エネルギー汎関数の凸性）
\item
  非線形波動方程式のソリトン解としての厳密な定式化
\item
  物理的時空との具体的な対応関係
\end{itemize}

本論文で「形不変定常波」と呼ぶ対象は、力学的含意を持つ「ソリトン」とは区別される。「形不変」とは平行移動に対して振幅分布
\(|\psi|\)
が保存されるという幾何学的性質のみを指し、「定常」とは節構造が球面族と整合する配置の性質のみを指す。

\begin{center}\rule{0.5\linewidth}{0.5pt}\end{center}

\subsection{\texorpdfstring{§2
閉弦上の形不変定常波（\(S^1\)）}{§2 閉弦上の形不変定常波（S\^{}1）}}\label{ux9589ux5f26ux4e0aux306eux5f62ux4e0dux5909ux5b9aux5e38ux6ce2s1}

\subsubsection{§2.1
パラメータ記述}\label{ux30d1ux30e9ux30e1ux30fcux30bfux8a18ux8ff0}

周長 \(L\)
の閉じた弦（\(S^1\)）上で、形不変定常波は以下の4パラメータで過不足なく記述される:

\begin{itemize}
\tightlist
\item
  \(\theta_0\)：最大波高位置の位置位相
\item
  \(\lambda_{\text{base}}\)：基底モード波長
\item
  \(A\)：最大波高の振幅
\item
  \(N\)：フーリエ展開のモード数（打ち切り次数）
\end{itemize}

ここで \(N\) は波形の「丸み／鋭さ」を制御するパラメータであり:

\begin{itemize}
\tightlist
\item
  \(N = 1\)：基底モードのみ（純粋な正弦波）
\item
  \(N\) 中間：奇数次高調波が加わった梯形的波形
\item
  \(N \to \infty\)：完全な矩形波（完全局在パルス）
\end{itemize}

\subsubsection{§2.2 定常波条件}\label{ux5b9aux5e38ux6ce2ux6761ux4ef6}

定常波が存在するための必要十分条件は:

\[L = k \cdot \lambda_{\text{base}} \quad (k \in \mathbb{Z}_{>0})\]

すなわち、閉弦の周長が基底モード波長の正整数倍であること。この条件下で、許される波長は自動的に
\(\lambda_n = L/n\)（\(n = 1, 2, 3, \ldots\)）に離散化される。

\subsubsection{§2.3
境界条件の統一}\label{ux5883ux754cux6761ux4ef6ux306eux7d71ux4e00}

位相波動方程式は弦が閉じていることを要求するが、閉じ方は指定しない。自由端反射（位相シフト
\(0\)）と固定端反射（位相シフト \(\pi\)）の違いは、鏡像法（method of
images）により長さ \(2L\)
の閉弦の偶パリティ/奇パリティ部分空間として統一的に扱える。境界条件の違いは基底モードの位相選択（\(0\)
or \(\pi\)）に吸収される。

\begin{center}\rule{0.5\linewidth}{0.5pt}\end{center}

\subsection{\texorpdfstring{§3 位相波動方程式と曲率 \(R\)
の直交性}{§3 位相波動方程式と曲率 R の直交性}}\label{ux4f4dux76f8ux6ce2ux52d5ux65b9ux7a0bux5f0fux3068ux66f2ux7387-r-ux306eux76f4ux4ea4ux6027}

\subsubsection{\texorpdfstring{§3.1 \(R\)
非依存性}{§3.1 R 非依存性}}\label{r-ux975eux4f9dux5b58ux6027}

位相波動方程式は局所的な微分方程式であり、空間曲率 \(R\)（あるいは
\(L = 2\pi R\)）はどこにも現れない。\(R\)
が関与するのは、「どの波長が定常波として成立するか」を選別する\textbf{大域的境界条件のレベルのみ}である:

\begin{itemize}
\tightlist
\item
  \(\lambda = L/n\)（\(n\) 整数）：一周で位相が \(2\pi n\)
  シフトし定常波が成立
\item
  \(\lambda > L\)：一周で位相が \(2\pi\)
  未満しかシフトせず定常波にならない
\item
  \(\lambda\) が \(L\)
  と無理数比：何周しても位相が元に戻らず、位相空間を準周期的に埋めていく
\end{itemize}

いずれの場合も波動方程式自体は同一であり、\(R\)
が小さくても大きくても波は同じ方程式に従う。

\subsubsection{\texorpdfstring{§3.2 \(R\)
の事後的定義}{§3.2 R の事後的定義}}\label{r-ux306eux4e8bux5f8cux7684ux5b9aux7fa9}

この構造から、\(R\)
は波動方程式の中にある量ではなく、\textbf{定常解の集合を眺めたときに事後的に読み取れる幾何量}として定義される:

\[R = \frac{\lambda_{\text{base}}}{2\pi}\]

すなわち \(R\) は「定常波として許容される最長波長の
\(1/(2\pi)\)」であり、先験的パラメータではなく導出量である。

\subsubsection{§3.3
構造的直交性}\label{ux69cbux9020ux7684ux76f4ux4ea4ux6027}

\textbf{これはユークリッド近似ではない。}
位相波動方程式と空間曲率は本質的に直交した演算であり、ユークリッド座標での計算が厳密に正しい。曲率
\(R\)
が関与するのは大域的な境界条件（閉じた空間の周長と波長の整合性）のみであり、波動方程式の局所的な解には一切影響しない。

\begin{center}\rule{0.5\linewidth}{0.5pt}\end{center}

\subsection{\texorpdfstring{§4 球面 \(S^2\)
への拡張と対蹠点収束}{§4 球面 S\^{}2 への拡張と対蹠点収束}}\label{ux7403ux9762-s2-ux3078ux306eux62e1ux5f35ux3068ux5bfeux8e60ux70b9ux53ceux675f}

\subsubsection{§4.1
パラメータ構造の保存}\label{ux30d1ux30e9ux30e1ux30fcux30bfux69cbux9020ux306eux4fddux5b58}

球面 \(S^2\) 上の形不変定常波は、\(S^1\)
と本質的に同じパラメータ構造を持つ。\(S^2\) の \(\mathrm{SO}(3)\)
対称性により、どの大円を選んでも定常波条件は同一の形:

\[2\pi R = n \cdot \lambda_{\text{base}}\]

パラメータは位置位相の次元だけ変化して
\((\theta_0, \lambda_{\text{base}}, A, N)\)
の4つで尽きる。波束の位置は完全に自由なパラメータであり、内部観測からは区別不可能である。

\subsubsection{§4.2
対蹠点収束現象}\label{ux5bfeux8e60ux70b9ux53ceux675fux73feux8c61}

\(S^2\) 上で大円周長が \(\lambda\)
の整数倍でない場合、一点に局在した波束は以下の運動を行う:

\begin{enumerate}
\def\labelenumi{\arabic{enumi}.}
\tightlist
\item
  初期：一点に局在した山
\item
  リング状（\(S^1\) 波面）に拡大
\item
  赤道付近でリング最大
\item
  反対側の対蹠点に向けて収縮
\item
  対蹠点で再び一点に収束
\item
  繰り返し
\end{enumerate}

\subsubsection{§4.3
幾何学的分散}\label{ux5e7eux4f55ux5b66ux7684ux5206ux6563}

\(S^2\) の固有振動数は \(\omega_l = c\sqrt{l(l+1)}/R\)
であり、\textbf{非等間隔}に並ぶ。これにより波束は厳密には周期的ではなく準周期的となり、サイクルを重ねるごとに波束の形が崩れていく。

\begin{center}\rule{0.5\linewidth}{0.5pt}\end{center}

\subsection{\texorpdfstring{§5 保存量
\(\int |u|^2 \, dV\)}{§5 保存量 \textbackslash int \textbar u\textbar\^{}2 \textbackslash, dV}}\label{ux4fddux5b58ux91cf-int-u2-dv}

\subsubsection{§5.1
瞬間波高と保存量の区別}\label{ux77acux9593ux6ce2ux9ad8ux3068ux4fddux5b58ux91cfux306eux533aux5225}

完全無損失の線形波動方程式において、瞬間的な波高 \(u(x,t)\)
そのものは局所的にも大域的にも「安定量」ではない。保存されるのは以下の非局所的積分量である:

\[\int_{S^n} |u|^2 \, dV = \sum_{l,m} |a_{lm}|^2 = \text{const}\]

波束の「同一性」は瞬間的な形状にではなく、この全空間にわたる2乗積分量に宿る。

\subsubsection{§5.2
シュレディンガー方程式との構造的同型性}\label{ux30b7ux30e5ux30ecux30c7ux30a3ux30f3ux30acux30fcux65b9ux7a0bux5f0fux3068ux306eux69cbux9020ux7684ux540cux578bux6027}

この保存則は量子力学のシュレディンガー方程式における
\(\int |\psi|^2 \, dV = 1\) と構造的に同型である:

{\def\LTcaptype{none} % do not increment counter
\begin{longtable}[]{@{}lll@{}}
\toprule\noalign{}
& 閉球面上の形不変定常波 & シュレディンガー方程式 \\
\midrule\noalign{}
\endhead
\bottomrule\noalign{}
\endlastfoot
場 & \(u(x,t)\)（実数の波高） & \(\psi(x,t)\)（複素振幅） \\
保存密度 & \(u^2\) & \(|\psi|^2\)（確率密度） \\
保存量 & \(\int u^2 \, dA\) & \(\int |\psi|^2 \, dV = 1\) \\
展開基底 & 球面調和関数 \(Y_{lm}\) & ハミルトニアン固有関数 \\
\end{longtable}
}

この同型性は偶然ではなく、「直交基底で展開できる線形方程式」という共通の数学的骨格から出る性質であり、古典か量子かは本質的ではない。

\begin{center}\rule{0.5\linewidth}{0.5pt}\end{center}

\subsection{\texorpdfstring{§6 \(S^3\) における Huygens
の原理の成立}{§6 S\^{}3 における Huygens の原理の成立}}\label{s3-ux306bux304aux3051ux308b-huygens-ux306eux539fux7406ux306eux6210ux7acb}

\subsubsection{§6.1
奇数次元と偶数次元の質的差異}\label{ux5947ux6570ux6b21ux5143ux3068ux5076ux6570ux6b21ux5143ux306eux8ceaux7684ux5deeux7570}

波動方程式の伝播性質は偶数空間次元と奇数空間次元で本質的に異なる（Hadamard
の結果）:

\begin{itemize}
\tightlist
\item
  \textbf{奇数次元}（\(S^1, S^3, S^5, \ldots\)）：厳密な Huygens
  の原理が成立。点源パルスは鋭い波面のみとして伝播し、後方に尾を引かない。
\item
  \textbf{偶数次元}（\(S^2, S^4, \ldots\)）：Huygens
  の原理が破れる。波面通過後も「尾（wake）」が残る。
\end{itemize}

\subsubsection{\texorpdfstring{§6.2 \(S^3\)
上の形不変定常波}{§6.2 S\^{}3 上の形不変定常波}}\label{s3-ux4e0aux306eux5f62ux4e0dux5909ux5b9aux5e38ux6ce2}

\(S^3\)（4次元空間内の完全対称な超球の超球殻）上では:

\begin{itemize}
\tightlist
\item
  保存量：\(\int |u|^2 \, dV\)（\(V = 2\pi^2 R^3\)）
\item
  定常波条件：\(2\pi R = n \cdot \lambda_{\text{base}}\)（\(\mathrm{SO}(4)\)
  対称性）
\item
  パラメータ：\((\theta_0, \phi_0, \psi_0, \lambda_{\text{base}}, A, N)\)
\end{itemize}

Huygens の原理により、一点に局在した波束は\textbf{鋭い球殻（\(S^2\)
状の波面）としてのみ伝播し}、後方に何も残さない。対蹠点で厳密に一点に収束して再び山が復元し、再び球殻として展開する。\(S^2\)
の残光を伴う準周期運動と比較して、\(S^3\)
ではより\textbf{クリーンな復元}が実現される。

\subsubsection{\texorpdfstring{§6.3 \(S^4\) での Huygens
の原理の再破}{§6.3 S\^{}4 での Huygens の原理の再破}}\label{s4-ux3067ux306e-huygens-ux306eux539fux7406ux306eux518dux7834}

\(S^4\) では Huygens の原理が再び破れ、\(S^2\)
と同様に波面の後方に残光が残る。\(S^3\)
で実現されたクリーンな復元が失われる。

\subsubsection{§6.4
次元と安定性の交代パターン}\label{ux6b21ux5143ux3068ux5b89ux5b9aux6027ux306eux4ea4ux4ee3ux30d1ux30bfux30fcux30f3}

{\def\LTcaptype{none} % do not increment counter
\begin{longtable}[]{@{}lll@{}}
\toprule\noalign{}
次元 & Huygens の原理 & 波束の振る舞い \\
\midrule\noalign{}
\endhead
\bottomrule\noalign{}
\endlastfoot
\(S^1\) & （1次元）厳密 & 位相回転のみ、完全周期 \\
\(S^2\) & 破れる & 対蹠点収束＋残光、準周期 \\
\(S^3\) & \textbf{成立} & 鋭い球殻伝播、クリーンな復元 \\
\(S^4\) & 破れる & \(S^2\) の高次元版、残光あり \\
\(S^5\) & \textbf{成立} & \(S^3\) の高次元版、クリーン \\
\end{longtable}
}

\begin{center}\rule{0.5\linewidth}{0.5pt}\end{center}

\subsection{§7
形不変定常波の安定存在条件}\label{ux5f62ux4e0dux5909ux5b9aux5e38ux6ce2ux306eux5b89ux5b9aux5b58ux5728ux6761ux4ef6}

\subsubsection{§7.1
偶数次元における波束同一性の喪失}\label{ux5076ux6570ux6b21ux5143ux306bux304aux3051ux308bux6ce2ux675fux540cux4e00ux6027ux306eux55aaux5931}

偶数次元（\(S^2, S^4\)）で起こる現象を整理する:

\begin{itemize}
\tightlist
\item
  保存量 \(\int |u|^2 \, dV\) は厳密保存される
\item
  各モード振幅 \(|a_{lm}|\) も個別に厳密保存される
\item
  しかし\textbf{局在した波束としての同一性}は保たれない
\end{itemize}

固有振動数 \(\omega_l\)
が無理数比で並ぶため、各モードの位相が独立にドリフトし、位相空間上で準周期運動する。長時間平均すると球面全体にわたる分布が典型状態となり、時折位相が偶然揃って局在が部分的に復活する（Poincaré
回帰）。これは KAM 理論的な準周期力学系の描像に一致する。

\subsubsection{§7.2
安定な形不変定常波の必要十分条件}\label{ux5b89ux5b9aux306aux5f62ux4e0dux5909ux5b9aux5e38ux6ce2ux306eux5fc5ux8981ux5341ux5206ux6761ux4ef6}

形不変定常波が「波束としての同一性を保って安定に存在できる」ための条件は:

\begin{enumerate}
\def\labelenumi{\arabic{enumi}.}
\tightlist
\item
  \textbf{Huygens の原理が成立すること}：波面が尾を引かない
\item
  \textbf{固有振動数が整数比に近いこと}：位相揃いが再生する
\end{enumerate}

これが両立するのは\textbf{奇数次元球面のみ}である。

\subsubsection{§7.3 最小性}\label{ux6700ux5c0fux6027}

\(S^1\) は Huygens
の原理を自明に満たすが、1次元では位相回転のみの退化した振る舞いしか示さない。非自明な波束の形成・復元・対蹠点収束を伴う最小の次元は
\textbf{\(S^3\)} である。

\textbf{命題}: 閉じた \(n\)
次元球面上で形不変定常波が非自明かつ安定に存在できる最小次元は \(n = 3\)
である。

\begin{center}\rule{0.5\linewidth}{0.5pt}\end{center}

\subsection{§8 Frobenius
の定理による独立論証}\label{frobenius-ux306eux5b9aux7406ux306bux3088ux308bux72ecux7acbux8ad6ux8a3c}

\subsubsection{§8.1
結合的除算代数}\label{ux7d50ux5408ux7684ux9664ux7b97ux4ee3ux6570}

Frobenius
の定理（1877）により、実数体上の有限次元結合的除算代数は以下の3つに限られる:

\begin{itemize}
\tightlist
\item
  \(\mathbb{R}\)（実数、1次元）
\item
  \(\mathbb{C}\)（複素数、2次元）
\item
  \(\mathbb{H}\)（四元数、4次元）
\end{itemize}

3次元には自然な除算代数が存在しない。代数構造（加減乗除・ノルム）を保ったまま拡張すると、3次元を飛ばして4次元に到達する。

\subsubsection{\texorpdfstring{§8.2 四元数の \(1+3\)
分解}{§8.2 四元数の 1+3 分解}}\label{ux56dbux5143ux6570ux306e-13-ux5206ux89e3}

四元数 \(q = a + bi + cj + dk\) は代数構造内在的に \(1\)
次元（スカラー部 \(a\)）\(+ 3\) 次元（ベクトル部
\(bi + cj + dk\)）に分解される。この分解は恣意的ではなく、四元数の共役、ノルム、乗法が要求する自然な構造である。

\subsubsection{§8.3
上限：八元数の非結合性}\label{ux4e0aux9650ux516bux5143ux6570ux306eux975eux7d50ux5408ux6027}

8次元の八元数 \(\mathbb{O}\) に拡張すると結合則が破れる（Hurwitz
の定理）。結合則を失うと波の合成・相互作用の記述が困難になるため、4次元が結合的除算代数としての最大次元である。

\begin{center}\rule{0.5\linewidth}{0.5pt}\end{center}

\subsection{§9
二重論証の収束}\label{ux4e8cux91cdux8ad6ux8a3cux306eux53ceux675f}

\subsubsection{§9.1
連続側と離散側の独立論証}\label{ux9023ux7d9aux5074ux3068ux96e2ux6563ux5074ux306eux72ecux7acbux8ad6ux8a3c}

{\def\LTcaptype{none} % do not increment counter
\begin{longtable}[]{@{}
  >{\raggedright\arraybackslash}p{(\linewidth - 4\tabcolsep) * \real{0.3333}}
  >{\raggedright\arraybackslash}p{(\linewidth - 4\tabcolsep) * \real{0.3333}}
  >{\raggedright\arraybackslash}p{(\linewidth - 4\tabcolsep) * \real{0.3333}}@{}}
\toprule\noalign{}
\begin{minipage}[b]{\linewidth}\raggedright
側面
\end{minipage} & \begin{minipage}[b]{\linewidth}\raggedright
結論
\end{minipage} & \begin{minipage}[b]{\linewidth}\raggedright
根拠
\end{minipage} \\
\midrule\noalign{}
\endhead
\bottomrule\noalign{}
\endlastfoot
連続波動（\(S^n\) 上の位相波動方程式） & 3次元空間が選好される &
奇数次元の Huygens の原理 \\
離散代数（格子構造） & 4次元が必然化される & Frobenius の定理 \\
統合 & \textbf{3 + 1 = 4} & 四元数の \(1+3\) 分解 \\
\end{longtable}
}

\subsubsection{§9.2 統合的命題}\label{ux7d71ux5408ux7684ux547dux984c}

連続側の「奇数次元で安定な形不変定常波」と離散側の「Frobenius
の定理」は、一見独立に見えて、\textbf{安定な形不変定常波が存在できる最小の時空}という単一の要請の二つの側面である:

\begin{enumerate}
\def\labelenumi{\arabic{enumi}.}
\tightlist
\item
  \textbf{空間構造}：波束の同一性を幾何学的に保護する空間 → 奇数次元球面
  → 最小非自明次元は \(S^3\)（3次元）
\item
  \textbf{代数構造}：波束の相互作用が同じカテゴリ内で閉じる →
  結合的除算代数 → 最小非自明次元は \(\mathbb{H}\)（4次元）
\item
  \textbf{統合}：四元数の \(1+3\) 分解により、3次元空間 \(+\)
  1次元時間軸として統一
\end{enumerate}

\textbf{定理（二重論証の収束）}:
形不変定常波が安定に存在しかつ代数的に閉じた相互作用を持つ最小の閉空間は
\(S^3\) であり、これを結合的除算代数内で記述する最小の構造は四元数
\(\mathbb{H}\) の \(1+3\) 分解、すなわち \(3+1\) 時空構造である。

\begin{center}\rule{0.5\linewidth}{0.5pt}\end{center}

\subsection{§10
位相波動方程式と空間曲率の直交性}\label{ux4f4dux76f8ux6ce2ux52d5ux65b9ux7a0bux5f0fux3068ux7a7aux9593ux66f2ux7387ux306eux76f4ux4ea4ux6027}

§3 で示した位相波動方程式と曲率 \(R\)
の直交性は、高次元でも保存される。\(S^1 \to S^2 \to S^3 \to S^4\)
のいずれにおいても:

\begin{itemize}
\tightlist
\item
  波動方程式は \(R\) に依存しない普遍的な方程式
\item
  \(R\) は定常波条件を通じて事後的に定まる幾何量
\item
  ユークリッド座標での計算が厳密に正しい（近似ではない）
\end{itemize}

この直交性は、形不変定常波が空間の歪みに無頓着であることを意味する。波は測地線に沿って伝播するが、空間が歪んでいるか否かは波動方程式の局所的な解には影響しない。

\begin{center}\rule{0.5\linewidth}{0.5pt}\end{center}

\subsection{§11 まとめ}\label{ux307eux3068ux3081}

本論文の主要な結果を要約する:

\begin{enumerate}
\def\labelenumi{\arabic{enumi}.}
\item
  \textbf{\(R\) 非依存性}:
  位相波動方程式と空間曲率は本質的に直交しており、\(R\)
  は定常解集合から事後的に読み取れる幾何量である。
\item
  \textbf{パラメータ記述の普遍性}: 形不変定常波のパラメータ
  \((\theta_0, \lambda_{\text{base}}, A, N)\) は \(S^1\) から \(S^n\)
  まで、位置位相の次元のみが変化する同一の骨格で記述される。
\item
  \textbf{保存量}: 波束の同一性を担うのは瞬間的な波高 \(u\)
  ではなく、非局所的な2乗積分 \(\int |u|^2 \, dV\) である。
\item
  \textbf{Huygens の原理による次元選択}:
  形不変定常波が安定に存在できるのは奇数次元球面に限られ、最小の非自明次元は
  \(S^3\)（3次元）である。
\item
  \textbf{二重論証の収束}: 連続側（Huygens の原理）と離散側（Frobenius
  の定理）が独立に同じ \(3+1\) 時空構造を必然化する。
\end{enumerate}

\begin{center}\rule{0.5\linewidth}{0.5pt}\end{center}

\subsection{参考文献}\label{ux53c2ux8003ux6587ux732e}

\begin{itemize}
\tightlist
\item
  {[}P1{]} 木原範昭, 「超球面間の中心投影」, 2025.
\item
  {[}P2{]} 木原範昭, 「超球殻間中心投影の正則性」, DOI:
  10.5281/zenodo.19692192, 2026.
\item
  {[}W8{]} 木原範昭, 「6次元超直方体と粒子模型」, 2025.
\item
  Hadamard, J., \emph{Lectures on Cauchy's Problem in Linear Partial
  Differential Equations}, 1923.
\item
  Frobenius, G., ``Über lineare Substitutionen und bilineare Formen'',
  \emph{J. reine angew. Math.}, \textbf{84}, 1--63, 1877.
\item
  Hurwitz, A., ``Über die Composition der quadratischen Formen von
  beliebig vielen Variablen'', \emph{Nachr. Ges. Wiss. Göttingen},
  309--316, 1898.
\end{itemize}
\end{document}
